\documentclass[10pt,aspectratio=169]{beamer}
\usepackage{poly}

% ============================================================================
%  PREÂMBULO  ->  colocar no main.tex, antes de \begin{document}
% ============================================================================
\usepackage{pgfplots}
\pgfplotsset{compat=1.18}
\usetikzlibrary{arrows.meta,calc}

\setbeamertemplate{caption}[numbered]
\setbeamertemplate{bibliography item}[text]

\usepackage[num]{abntex2cite}
\citebrackets{[}{]}
\hypersetup{hidelinks}


% ------------------------------------------------------------------ ajustes
\definecolor{sinalcor}{HTML}{1F4E79}
\newcommand{\sinal}{sin(deg(2*pi*0.75*x + 0.4))}   % x(t)
\newcommand{\Ts}{0.11}                              % periodo de amostragem
\newcommand{\Tau}{0.0367}                           % largura do pulso
\newcommand{\Nam}{9}                                % ultima amostra

\pgfplotsset{
    pamaxis/.style={
        width=\textwidth, height=4.2cm,
        axis lines=middle,
        axis line style={-{Stealth[length=5pt,width=4pt]}},
        xmin=-0.04, xmax=1.16, ymin=-1.15, ymax=1.15,
        xtick=\empty, ytick=\empty,
        clip=false, samples=201, domain=0:1,
        xlabel=$t$,
        xlabel style={at={(axis description cs:1,0.5)}, anchor=west},
        ylabel style={at={(axis description cs:0.033,1)}, anchor=south east},
    },
}

\newcommand{\marcaTs}{%
    \draw[gray, {Stealth[length=4pt]}-{Stealth[length=4pt]}]
        (axis cs:0.11,-0.55) -- node[below=-2pt, font=\small, black] {$T_s$}
        (axis cs:0.22,-0.55);}

% ============================================================================
%  Slides do método topo-plano (topicos/flattop.tex)
%  Todas as macros usam o prefixo \ft para não conflitar com \Ts e \Tau acima.
% ============================================================================
\usepgfplotslibrary{groupplots}

% ---------------- Parâmetros (altere aqui) ----------------
\newcommand{\ftfzero}{1.6}  % frequência de x(t) após o filtro anti-aliasing [Hz]
\newcommand{\ftfs}{8}       % frequência de amostragem f_s [Hz] (parâmetro fixo do estudo)
\newcommand{\ftduty}{0.5}   % razão tau/T_s (0 < duty <= 1)
\newcommand{\fttmax}{1.25}  % duração exibida [s] (2 períodos de 1,6 Hz)
\newcommand{\ftpico}{0.5}   % pico onde o atraso é indicado em (c) e (d)
% -----------------------------------------------------------

\pgfmathsetmacro\ftTs{1/\ftfs}
\pgfmathsetmacro\fttau{\ftduty*\ftTs}
\pgfmathsetmacro\ftA{sin(deg(pi*\ftfzero*\fttau))/(pi*\ftfzero*\fttau)} % fator de abertura
\pgfmathsetmacro\ftpicoatr{\ftpico+\fttau/2}                           % pico atrasado
\pgfmathtruncatemacro\ftN{\fttmax*\ftfs}

% Constrói os pulsos e as amostras
\def\ftpulsos{}\def\ftamostras{}\def\ftamospos{}\def\ftamosneg{}\def\ftimpt{}
\foreach \n in {0,...,\ftN}{%
  \pgfmathsetmacro\fta{\n*\ftTs}%
  \pgfmathparse{\fta<\fttmax ? 1 : 0}%
  \ifnum\pgfmathresult=1
    \pgfmathsetmacro\ftb{min(\fta+\fttau,\fttmax)}%
    \pgfmathsetmacro\ftv{cos(deg(2*pi*\ftfzero*\fta))}%
    \xdef\ftpulsos{\ftpulsos (\fta,0) (\fta,\ftv) (\ftb,\ftv) (\ftb,0)}%
    \xdef\ftamostras{\ftamostras (\fta,\ftv)}%
    \xdef\ftimpt{\ftimpt (\fta,1)}%
    \pgfmathparse{\ftv>=0 ? 1 : 0}%
    \ifnum\pgfmathresult=1 \xdef\ftamospos{\ftamospos (\fta,\ftv)}%
    \else \xdef\ftamosneg{\ftamosneg (\fta,\ftv)}\fi
  \fi
}

\definecolor{ftazul}{RGB}{42,120,214}
\definecolor{ftlaranja}{RGB}{220,95,40}
\definecolor{ftroxo}{RGB}{98,80,214}
\definecolor{ftverde}{RGB}{0,131,0}

\pgfplotsset{
  ftaxis/.style={
    scale only axis, width=0.19\textwidth, height=0.40\textheight,
    xmin=0, xmax=\fttmax, ymin=-1.2, ymax=1.3,
    xtick={0,0.5,...,\fttmax}, ytick={-1,-0.5,0,0.5,1},
    xticklabel style={/pgf/number format/use comma},
    yticklabel style={/pgf/number format/use comma},
    xlabel={$t$ (s)}, grid=major, grid style={gray!15},
    tick label style={font=\tiny}, label style={font=\footnotesize},
    title style={font=\scriptsize, align=center},
    axis lines=left, enlarge x limits=false,
    extra y ticks={0}, extra y tick labels={},
    extra y tick style={grid style={black!50, thin}},
  }
}


% ============================================================================
%  Slide "Representação Gráfica no Domínio da Frequência" (topicos/flattop.tex)
%  Usa os parâmetros \ftfzero, \ftfs, \fttau e \ftA definidos acima.
%  Espectros normalizados por tau/T_s.
% ============================================================================
\newcommand{\ftfmax}{18}    % faixa de frequência exibida: -fmax a fmax [Hz]
\pgfmathsetmacro\ftfc{\ftfs/2}          % corte do passa-baixas ideal: f_s/2
\pgfmathsetmacro\ftnulo{1/\fttau}       % primeiro nulo do sinc: 1/tau

% Réplicas X(f - k f_s) ponderadas por |sinc(f tau)|/2
\def\ftreplicas{}
\foreach \k in {-5,...,5}{%
  \foreach \s in {-1,1}{%
    \pgfmathsetmacro\ftf{\s*\ftfzero+\k*\ftfs}%
    \pgfmathparse{abs(\ftf)<=\ftfmax ? 1 : 0}%
    \ifnum\pgfmathresult=1
      \pgfmathsetmacro\fth{0.5*abs(sin(deg(pi*\ftf*\fttau))/(pi*\ftf*\fttau))}%
      \xdef\ftreplicas{\ftreplicas (\ftf,\fth)}%
    \fi
  }%
}
\pgfmathsetmacro\ftAmeio{0.5*\ftA}
\newcommand{\ftstrut}{\rule[-1ex]{0pt}{3ex}}   % alinha as 2as linhas dos títulos

\pgfplotsset{
  ftfreqaxis/.style={
    ftaxis,
    xmin=-\ftfmax, xmax=\ftfmax, ymin=0, ymax=0.68,
    xtick={-10,0,10}, minor xtick={-15,-5,5,15},
    ytick={0,0.25,0.5},
    xlabel={$f$ (Hz)}, axis lines=left,
    extra y ticks={}, extra x ticks={},
    clip mode=individual,
  },
  ftimpulso/.style={ycomb, thick, mark=triangle*, mark size=1.6pt},
}

% ============================================================================
%  Slides de demonstração (tempo e frequência)
% ============================================================================
% Réplicas sem ponderação (amostragem ideal) e ponderadas por |sinc|
\def\ftrepideal{}\def\ftrepsinc{}\def\ftimpf{}
\foreach \k in {-5,...,5}{%
  \pgfmathsetmacro\ftf{\k*\ftfs}%
  \pgfmathparse{abs(\ftf)<=\ftfmax ? 1 : 0}%
  \ifnum\pgfmathresult=1 \xdef\ftimpf{\ftimpf (\ftf,1)}\fi
  \foreach \s in {-1,1}{%
    \pgfmathsetmacro\ftf{\s*\ftfzero+\k*\ftfs}%
    \pgfmathparse{abs(\ftf)<=\ftfmax ? 1 : 0}%
    \ifnum\pgfmathresult=1
      \pgfmathsetmacro\fth{abs(sin(deg(pi*\ftf*\fttau))/(pi*\ftf*\fttau))}%
      \xdef\ftrepideal{\ftrepideal (\ftf,1)}%
      \xdef\ftrepsinc{\ftrepsinc (\ftf,\fth)}%
    \fi
  }%
}

\pgfplotsset{
  ftdemot/.style={
    ftaxis, width=0.23\textwidth, height=0.22\textheight,
    title style={font=\footnotesize}, clip mode=individual,
  },
  ftdemof/.style={
    ftdemot,
    xmin=-\ftfmax, xmax=\ftfmax, ymin=0, ymax=1.3,
    xtick={-10,-5,0,5,10}, minor xtick={},
    ytick={1}, extra y ticks={},
    xlabel={$f$ (Hz)},
  },
  ftimpneg/.style={only marks, mark=triangle*, mark size=1.6pt, mark options={rotate=180}},
  ftimppos/.style={only marks, mark=triangle*, mark size=1.6pt},
}
% símbolo de operação entre dois gráficos de um groupplot chamado "g"
\newcommand{\ftop}[3]{\node[font=\Large] at ($(g c#1r1.east)!0.5!(g c#2r1.west)$) {$#3$};}

% ============================================================================
%  Numeração dos slides (rodapé à direita: atual/total)
% ============================================================================
\setbeamertemplate{footline}{%
  \hfill\usebeamerfont{footline}\color{gray}%
  \insertframenumber/\inserttotalframenumber\hspace{2ex}\vskip2pt%
}

% ============================================================================
%  Últimos slides (abertura, deslocamento, reconstrução)
% ============================================================================
% Overlays em elementos do pgfplots: \addplot[visible on=<2->] ...
\tikzset{
  invisible/.style={opacity=0, text opacity=0},
  visible on/.style={alt={#1{}{invisible}}},
  alt/.code args={<#1>#2#3}{\alt<#1>{\pgfkeysalso{#2}}{\pgfkeysalso{#3}}},
}
\newcommand{\ftverde}[1]{{\color{ftverde}#1}}
% Núcleo de interpolação normalizado: g(t)/tau = (1/tau) int_0^tau sinc(f_s(t - l)) dl
% (regra do ponto médio com 5 pontos; o deslocamento 1e-4 evita a divisão 0/0)
\pgfmathdeclarefunction{ftsincfs}{1}{%
  \pgfmathparse{sin(deg(pi*\ftfs*(#1)+1.234e-4))/(pi*\ftfs*(#1)+1.234e-4)}}
\pgfmathdeclarefunction{ftnucleo}{1}{%
  \pgfmathparse{0.2*(ftsincfs(#1-0.1*\fttau)+ftsincfs(#1-0.3*\fttau)+ftsincfs(#1-0.5*\fttau)
    +ftsincfs(#1-0.7*\fttau)+ftsincfs(#1-0.9*\fttau))}}
% Desenha os núcleos x(nT_s) g(t - nT_s)/tau (usado no slide de reconstrução no tempo)
\newcommand{\ftnucleos}{%
  \pgfplotsinvokeforeach{-3,...,13}{%
    \addplot[ftroxo!55, thin, domain=0:\fttmax, samples=150, visible on=<2->]
      {cos(deg(2*pi*\ftfzero*##1*\ftTs))*ftnucleo(x-##1*\ftTs)};
  }%
}

% ============================================================================
%  Amostragem IDEAL (topicos/ideal.tex) — prefixo \id
%  Reaproveita \ftfzero, \ftfs, \ftTs, \fttmax e os estilos ft*.
%  Diferença essencial: não há pulso h(t), logo as réplicas não são
%  ponderadas pelo sinc e não existem abertura nem atraso.
% ============================================================================
% Réplicas de amostragem ideal com altura 1/2 (mesma escala dos slides ft)
\def\idrepmeio{}
\foreach \k in {-5,...,5}{%
  \foreach \s in {-1,1}{%
    \pgfmathsetmacro\idf{\s*\ftfzero+\k*\ftfs}%
    \pgfmathparse{abs(\idf)<=\ftfmax ? 1 : 0}%
    \ifnum\pgfmathresult=1 \xdef\idrepmeio{\idrepmeio (\idf,0.5)}\fi
  }%
}
% Núcleos de interpolação x(nT_s) sinc(f_s(t - nT_s))
\newcommand{\idnucleos}{%
  \pgfplotsinvokeforeach{-4,...,14}{%
    \addplot[ftroxo!55, thin, domain=0:\fttmax, samples=200, visible on=<2->]
      {cos(deg(2*pi*\ftfzero*##1*\ftTs))*ftsincfs(x-##1*\ftTs)};
  }%
}
% Cores das etapas na amostragem ideal
\newcommand{\idamos}[1]{{\color{ftazul}#1}}            % trem de impulsos
\newcommand{\idrec}[1]{{\color{ftverde!70!black}#1}}   % reconstrução

% ============================================================
% Metadata
% ============================================================

\title{Análise Comparativa do Processo de Amostragem para Sinais Analógicos}
\subtitle{TEC513 - MI - Processamento Digital de Sinais}
\author{Douglas Oliveira, Júlia Carneiro, Yasmin Cordeiro}
\institute[Engenharia da Computação]{Universidade Estadual de Feira de Santana}
\date{\today}


% =============================================================
% Main Body
% =============================================================

\begin{document}
\maketitle % generate the title slide

% Sinais Digitais e Analógicos

% Introdução - requisitos

% Amostragem
% 	O Teorema da Amostragem (Nyquist-Shannon)
% 	Aliasing (Subamostragem)
% 	Filtro Anti-Aliasing 

% Parâmetros fixos

% Método Ideal:
% - Descrição Analítica
% - Simulação de amostragem
% - Filtro de reconstrução
% - Simulação de reconstrução

% Método Natural:
% - Descrição Analítica
% - Simulação de amostragem
% - Filtro de reconstrução
% - Simulação de reconstrução

% Método Topo-Plano:
% - Descrição Analítica
% - Simulação de amostragem
% - Filtro de reconstrução
% - Simulação de reconstrução

% Cenários de Teste:
% - Casos de aliasing:
% 	- Variando f de corte e f de amostragem
% - Tau:
% 	- Aumentar/diminuir e explicar as consequências
% - Vetor de frequências:
% 	- Aumentar e diminuir o vetor para demostrar a sobreposições de sincs

% Resultados e Discussões:
% - Comparação entre métodos
% - Comparar em tempo e frequência
% 
% Scripts não precisam ser mostrados na apresentação, mas vão ser entregues e podem ser citados em questionamentos. GRAFICO DE MAGNITUDE E FASER PARA TOPO PLANO


\section{Introdução}

\begin{frame}{Contextualização}

    A empresa \textit{SigmaDSP Inc.} enfrenta o desafio de compreender minuciosamente o \textbf{teorema de amostragem} de sinais analógicos de tempo contínuo. Para isso sabe-se que três são as formas ou
métodos da amostragem de sinais analógicos de tempo contínuo:
    \begin{itemize}
        \item Ideal
        \item Natural 
        \item Instantânea de topo plano \textit{(flat-top)}.
    \end{itemize}
    Partindo desse contexto, foram solicitadas as simulações dos \textbf{métodos de amostragem} para uma senoide e os respectivos filtros de reconstrução. 
    
\end{frame}

\begin{frame}{Requisitos}

	Os requisitos para a resolução do problema eram:
	\begin{itemize}
		\item Simulação dos métodos de amostragem com uma senoide 
            \begin{itemize}
                \item \textit{Com uma segunda senoide sendo utilizada como impulso de frequência acima da máxima permitida pelo filtro.}
            \end{itemize}
        \item Frequência de amostragem ajustável 
            \begin{itemize}
                \item Para conseguir demonstrar o fenômeno do \textit{\textbf{aliasing}}
            \end{itemize}
        \item Filtragem de reconstrução
        \item Descrição e comparação analítica de todos os métodos e suas reconstruções. 
        
		
	\end{itemize}
\end{frame}
% ============================================================================
%  Fundamentação Teórica
%  Abertura da seção: o que os três métodos têm em comum e o que os diferencia.
%  Em seguida, cada método é desenvolvido na ordem:
%     (1) descrição analítica geral  (2) aplicação aos parâmetros fixos
%     (3) reconstrução               (4) gancho para o método seguinte
% ============================================================================

\section{Fundamentação Teórica}

\begin{frame}{Amostragem}{Estrutura Comum aos Três Métodos}
    \small
    Amostrar é registrar o sinal contínuo apenas nos instantes $t = nT_s$. Os três
    métodos partem do mesmo trem de impulsos e diferem somente no \textbf{pulso} que
    carrega cada amostra:
    \begin{equation*}
        x_{am}(t) = \underbrace{\left[\, x_f(t) \cdot \sum_{n=-\infty}^{\infty} \delta(t - nT_s) \right]}_{\text{comum aos três métodos}}
        \;\longrightarrow\; \text{pulso de cada método}
    \end{equation*}

    \begin{center}
    \begin{tabular}{lll}
        \hline
        \textbf{Método} & \textbf{Amostra transportada por} & \textbf{Consequência} \\
        \hline
        Ideal & impulso $\delta(t - nT_s)$ & recuperação exata \\
        Natural & pulso que \textit{acompanha} $x_f(t)$ & espectro ponderado por $\mathrm{sinc}$ \\
        Topo-plano & pulso de topo constante $h(t)$ & $\mathrm{sinc}$ + atraso $\tau/2$ \\
        \hline
    \end{tabular}
    \end{center}

    \vspace{0.5em}
    Começamos pelo \textbf{ideal}: é o caso sem pulso, e portanto a referência com que
    os outros dois serão comparados.
\end{frame}

% ============================================================================
%  Método Ideal (amostragem por trem de impulsos)
%
%  Ordem dos slides:
%    A. Descrição analítica geral  (conceito, interpretação gráfica,
%       demonstração no tempo e na frequência)
%    B. Aplicação aos parâmetros fixos  (x_f(t) = cos(2 pi 2 t), f_s = 5 Hz)
%    C. Reconstrução  (filtro, frequência, tempo e verificação numérica)
%    D. Visão geral do processo + gancho para os métodos com pulso
%
%  Macros \ft... e \id... definidas no main.tex.
%  As figuras usam os próprios parâmetros fixos: f_0 = 2 Hz e f_s = 5 Hz.
% ============================================================================

% ----------------------------------------------------------------------------
%  A. DESCRIÇÃO ANALÍTICA GERAL
% ----------------------------------------------------------------------------
\begin{frame}{Método Ideal: Amostragem}{Análise no Domínio do Tempo}

    Multiplicamos o sinal $x_f(t)$ por um trem de impulsos deslocados de $T_s$.
    Cada impulso ``recolhe'' o valor instantâneo do sinal, sem qualquer retenção:
    a amostra existe apenas no instante $t = nT_s$.

    \vspace{1em}
    A equação para a modulação é dada por:
    \begin{equation*}
        x_\delta(t) = x_f(t) \cdot p(t),
        \qquad
        p(t) = \sum_{n=-\infty}^{\infty} \delta(t - nT_s)
    \end{equation*}

    Pela propriedade de peneiramento do impulso,
    $x_f(t)\,\delta(t - nT_s) = x_f(nT_s)\,\delta(t - nT_s)$:
    \begin{equation*}
        x_\delta(t) = \sum_{n=-\infty}^{\infty} x_f(nT_s)\, \delta(t - nT_s)
    \end{equation*}

    Não há pulso de retenção: o trem de impulsos é o único elemento do processo.

\end{frame}

\begin{frame}{Método Ideal: Amostragem}{Interpretação Gráfica}
\centering
\begin{minipage}[c][0.66\textheight][c]{0.92\textwidth}
\centering

% ------------------------------------------------------- (1) sinal original
\only<1>{%
\begin{tikzpicture}
    \begin{axis}[pamaxis, ylabel={$x_f(t)$}]
        \addplot[sinalcor, line width=1pt, smooth] {\sinal};
    \end{axis}
\end{tikzpicture}\par
Sinal analógico de banda base $x_f(t)$.}

% ------------------------------------------------------ (2) trem de impulsos
\only<2>{%
\begin{tikzpicture}
    \begin{axis}[pamaxis, ylabel={$p(t)$}]
        \addplot[draw=none] {0};
        \pgfplotsextra{%
            \foreach \n in {0,...,\Nam}{%
                \pgfmathsetmacro\xn{\n*\Ts}%
                \draw[sinalcor, line width=1pt, -{Stealth[length=5pt,width=4pt]}]
                    (axis cs:\xn,0) -- (axis cs:\xn,0.85);
                \draw[gray, line width=.4pt] (axis cs:\xn,-0.07) -- (axis cs:\xn,0.07);}}
        \marcaTs
    \end{axis}
\end{tikzpicture}\par
Trem de impulsos com período $T_s$, todos de área unitária.}

% ------------------------------------------- (3) produto (amostragem ideal)
\only<3>{%
\begin{tikzpicture}
    \begin{axis}[pamaxis, ylabel={$x_\delta(t)$}]
        \addplot[gray!50, dashed, line width=.7pt, smooth] {\sinal};
        \pgfplotsextra{%
            \foreach \n in {0,...,\Nam}{%
                \pgfmathsetmacro\xn{\n*\Ts}%
                \pgfmathsetmacro\yn{sin(deg(2*pi*0.75*\n*\Ts + 0.4))}%
                \draw[sinalcor, line width=1pt, -{Stealth[length=5pt,width=4pt]}]
                    (axis cs:\xn,0) -- (axis cs:\xn,\yn);}}
        \marcaTs
    \end{axis}
\end{tikzpicture}\par
Produto: impulsos de área $x_f(nT_s)$, nulos fora dos instantes de amostragem.}

\end{minipage}
\end{frame}

\begin{frame}{Método Ideal: Amostragem}{Demonstração no Domínio do Tempo}
\small
No domínio do tempo, o processo é apenas a \idamos{multiplicação do sinal por um trem
de impulsos}, sem etapa de retenção. Aplicando o peneiramento, cada impulso fica
ponderado pela sua amostra:
\begin{equation*}
  x_\delta(t) = x_f(t)\cdot \idamos{p(t)}
  = \boxed{\;\sum_{n=-\infty}^{\infty} x_f(nT_s)\, \delta(t - nT_s)\;},
  \qquad
  \idamos{p(t) = \sum_{n=-\infty}^{\infty} \delta(t - nT_s)}
\end{equation*}
\begin{center}
\begin{tikzpicture}
\begin{groupplot}[group style={group name=g, group size=3 by 1, horizontal sep=1.3cm}, ftdemot, height=0.18\textheight]
\nextgroupplot[title={(a) Sinal original\\[1pt]{\tiny\ftstrut $x_f(t)$}}]
  \addplot[ftazul!40!black, thick, domain=0:\fttmax, samples=300] {cos(deg(2*pi*\ftfzero*x))};
\nextgroupplot[title={(b) Trem de impulsos\\[1pt]{\tiny\ftstrut $p(t)$}}]
  \addplot[ycomb, ftazul, thick] coordinates {\ftimpt};
  \addplot[ftimppos, ftazul] coordinates {\ftimpt};
  \draw[gray, <->, >=stealth] (axis cs:\ftTs,1.18) -- node[above, font=\tiny, inner sep=1pt]{$T_s$} (axis cs:2*\ftTs,1.18);
\nextgroupplot[title={(c) Sinal amostrado\\[1pt]{\tiny\ftstrut $x_\delta(t)$}}]
  \addplot[gray!55, dashed, domain=0:\fttmax, samples=300] {cos(deg(2*pi*\ftfzero*x))};
  \addplot[ycomb, ftazul, thick] coordinates {\ftamostras};
  \addplot[ftimppos, ftazul] coordinates {\ftamospos};
  \addplot[ftimpneg, ftazul] coordinates {\ftamosneg};
\end{groupplot}
\ftop{1}{2}{\times}
\ftop{2}{3}{=}
\end{tikzpicture}
\end{center}
O sinal amostrado só existe nos instantes $t = nT_s$: entre as amostras não há informação.
\end{frame}

\begin{frame}{Método Ideal: Amostragem}{Demonstração no Domínio da Frequência}
\small
Aplicando a transformada de Fourier, o produto no tempo vira convolução na frequência.
Como a convolução com um impulso desloca o espectro, cada impulso de
$\idamos{P(f)}$ cria uma cópia inteira de $X_f(f)$:
\begin{equation*}
  X_\delta(f) = X_f(f) * \idamos{P(f)},
  \quad
  \idamos{P(f) = \frac{1}{T_s}\sum_{n=-\infty}^{\infty} \delta(f - nf_s)}
  \;\Longrightarrow\;
  \boxed{\; X_\delta(f) = f_s\sum_{n=-\infty}^{\infty} X_f(f - nf_s) \;}
\end{equation*}
\begin{center}
\begin{tikzpicture}
\begin{groupplot}[group style={group name=g, group size=3 by 1, horizontal sep=1.3cm}, ftdemof, height=0.18\textheight]
\nextgroupplot[title={(a) Espectro original\\[1pt]{\tiny\ftstrut $X_f(f)$}}, yticklabels={$\frac{1}{2}$}]
  \addplot[ftimpulso, ftazul!40!black] coordinates {(-\ftfzero,1) (\ftfzero,1)};
\nextgroupplot[title={(b) Trem de impulsos\\[1pt]{\tiny\ftstrut $P(f)$}}, yticklabels={$\frac{1}{T_s}$}]
  \addplot[ftimpulso, ftazul] coordinates {\ftimpf};
\nextgroupplot[title={(c) Sinal amostrado\\[1pt]{\tiny\ftstrut $X_\delta(f)$}}, yticklabels={$\frac{1}{2T_s}$}]
  \addplot[ftimpulso, ftazul] coordinates {\ftrepideal};
\end{groupplot}
\ftop{1}{2}{*}
\ftop{2}{3}{=}
\end{tikzpicture}
\end{center}
Todas as réplicas têm a mesma amplitude: não há envelope ponderando o espectro.
\end{frame}

\begin{frame}{Método Ideal: Amostragem}{Condição de Recuperação}
\small
A banda base ocupa $|f| \le f_{max}$ e a réplica vizinha começa em $f_s - f_{max}$.
Para que não haja sobreposição:
\begin{equation*}
  f_s - f_{max} > f_{max}
  \quad\Longrightarrow\quad
  \boxed{\; f_s > 2 f_{max} \;}
  \qquad
  \text{(banda de guarda } = f_s - 2f_{max})
\end{equation*}
\begin{itemize}
  \item Satisfeita a condição, as réplicas ficam separadas e a banda base pode ser
        isolada --- é o que torna a reconstrução possível.
  \item Violada a condição, as réplicas invadem a banda base: \textit{aliasing}, que
        nenhum filtro posterior desfaz.
  \item Por isso o filtro anti-aliasing atua \textbf{antes} da amostragem, limitando
        $f_{max}$ ao valor previsto no projeto.
\end{itemize}
\end{frame}

% ----------------------------------------------------------------------------
%  B. APLICAÇÃO AOS PARÂMETROS FIXOS
% ----------------------------------------------------------------------------
\begin{frame}{Método Ideal: Aplicação}{Parâmetros Fixos --- Domínio do Tempo}
    \small
    Com $x_f(t) = \cos(2\pi \cdot 2t)$, $f_s = 5$~Hz e $T_s = 0,2$~s:
    \begin{align*}
        x_\delta(t) &= x_f(t) \cdot \sum_{n=-\infty}^{\infty} \delta(t - nT_s)
        = \sum_{n=-\infty}^{\infty} x_f(nT_s)\, \delta(t - nT_s) \\
        &= \sum_{n=-\infty}^{\infty} \cos\!\left(2\pi \cdot 2 \cdot \frac{n}{5}\right) \delta\!\left(t - \frac{n}{5}\right)
        = \sum_{n=-\infty}^{\infty} \cos\!\left(\frac{4\pi n}{5}\right) \delta\!\left(t - \frac{n}{5}\right)
    \end{align*}

    \textbf{Amostras (o padrão se repete a cada 5 amostras, pois $f_0/f_s = 2/5$):}
    \[
        \begin{array}{c|ccccc}
            n & 0 & 1 & 2 & 3 & 4 \\ \hline
            x_f(nT_s) & 1 & -0{,}809 & 0{,}309 & 0{,}309 & -0{,}809
        \end{array}
    \]

    São $2{,}5$ amostras por período do sinal: o mínimo exigido por $f_s > 2f_{max}$.
\end{frame}

\begin{frame}{Método Ideal: Aplicação}{Parâmetros Fixos --- Domínio da Frequência}
    \small
    Substituindo $f_s = 5$~Hz na expressão do espectro amostrado:
    $$X_\delta(f) = f_s \cdot \sum_{k=-\infty}^{\infty} X_f(f - k f_s) = 5 \cdot \sum_{k=-\infty}^{\infty} X_f(f - 5k)$$

    \textbf{Análise das componentes:}
    $$
    \left.\begin{matrix}
    k=0: \pm 2\text{ Hz} \\
    k=1: 3 \text{ e } 7\text{ Hz} \\
    k=-1: -7 \text{ e } -3\text{ Hz} \\
    k=2: 8 \text{ e } 12\text{ Hz}
    \end{matrix}\right\} \text{componentes em } \pm 2, \pm 3, \pm 7, \pm 8, \pm 12, \dots
    $$

    Cada réplica tem peso $f_s \cdot \frac{1}{2} = 2{,}5$: \textbf{todas com a mesma
    amplitude}, pois não há pulso ponderando o espectro. A banda base ($2$~Hz) e a
    primeira réplica ($3$~Hz) estão separadas --- não há \textit{aliasing}.
\end{frame}

% ----------------------------------------------------------------------------
%  C. RECONSTRUÇÃO
% ----------------------------------------------------------------------------
\begin{frame}{Método Ideal: Reconstrução}{Filtro Passa-Baixas Ideal}
    \small
    Separadas as réplicas, basta descartar todas menos a banda base. O filtro de
    reconstrução é um passa-baixas ideal de corte $f_s/2$ e ganho $T_s$:
    \[
        H_{PB}(f) =
        \begin{cases}
            T_s, & |f| < \dfrac{f_s}{2}\\[4pt]
            0, & \text{caso contrário}
        \end{cases}
        \qquad\longrightarrow\qquad
        \text{com os parâmetros fixos: }
        \begin{cases}
            0{,}2, & |f| < 2{,}5\,\text{Hz}\\
            0, & \text{caso contrário}
        \end{cases}
    \]

    \textbf{Verificação da separação:}
    \[
        \underbrace{2\,\text{Hz}}_{\text{banda base}} < \underbrace{2{,}5\,\text{Hz}}_{\text{corte}} < \underbrace{3\,\text{Hz}}_{\text{1ª réplica}}
        \qquad \Longrightarrow \qquad
        \text{margem de } 0{,}5\,\text{Hz de cada lado}
    \]

    Essa margem é a banda de guarda de $1$~Hz obtida com $f_s = 5$~Hz: se
    $f_s \le 4$~Hz, a réplica em $f_s - 2$ entraria na banda base.
\end{frame}

\begin{frame}{Método Ideal: Reconstrução}{Demonstração no Domínio da Frequência}
\small
A reconstrução é o produto do espectro amostrado pelo passa-baixas. O ganho $T_s$
cancela exatamente o fator $f_s$ introduzido pela amostragem:
\begin{equation*}
  X_r(f) = X_\delta(f)\cdot \idrec{H_{PB}(f)} = T_s\cdot\frac{1}{T_s}X_f(f) = \boxed{\,X_f(f)\,}
  \qquad
  \text{(com os valores: } 0{,}2 \cdot 5 \cdot \tfrac{1}{2} = \tfrac{1}{2}\text{)}
\end{equation*}
\begin{center}
\begin{tikzpicture}
\begin{groupplot}[group style={group name=g, group size=3 by 1, horizontal sep=1.3cm}, ftdemof, height=0.18\textheight]
\nextgroupplot[title={(a) Espectro amostrado\\[1pt]{\tiny\ftstrut $X_\delta(f)$ e $H_{PB}(f)$}}, yticklabels={$\frac{1}{2T_s}$}]
  \draw[ftroxo, densely dashed, fill=ftroxo!8] (axis cs:-\ftfc,0) rectangle (axis cs:\ftfc,1.15);
  \addplot[ftimpulso, ftazul] coordinates {\ftrepideal};
\nextgroupplot[title={(b) Passa-baixas ideal\\[1pt]{\tiny\ftstrut $H_{PB}(f)/T_s$}}, yticklabels={$1$},
  visible on=<2->]
  \addplot[ftroxo, thick, fill=ftroxo!8, visible on=<2->] coordinates
    {(-\ftfc,0) (-\ftfc,1) (\ftfc,1) (\ftfc,0)};
\nextgroupplot[title={(c) Sinal reconstruído\\[1pt]{\tiny\ftstrut $X_r(f) = X_f(f)$}}, yticklabels={$\frac{1}{2}$}]
  \addplot[ftimpulso, ftverde, visible on=<3->] coordinates {(-\ftfzero,1) (\ftfzero,1)};
\end{groupplot}
\ftop{1}{2}{\times}
\ftop{2}{3}{=}
\end{tikzpicture}
\end{center}
Sem ponderação por $\mathrm{sinc}$ e sem termo de fase: o espectro recuperado é
\textbf{idêntico} ao original, dispensando equalizador.
\pause\pause
\end{frame}

\begin{frame}{Método Ideal: Reconstrução}{Demonstração no Domínio do Tempo}
\small
No tempo, a reconstrução é a convolução do sinal amostrado com a resposta ao impulso do
passa-baixas, $h_{PB}(t) = \mathrm{sinc}(f_s t)$. Cada impulso vira um $\mathrm{sinc}$
ponderado pela sua amostra:
\begin{equation*}
  x_r(t) = x_\delta(t) * \idrec{h_{PB}(t)}
  = \sum_{n=-\infty}^{\infty} x_f(nT_s)\, \idrec{\mathrm{sinc}\!\big(f_s(t - nT_s)\big)}
  = \boxed{\, x_f(t) \,}
\end{equation*}
A interpolação é exata porque $\mathrm{sinc}(f_s(t-nT_s))$ vale $1$ em $t = nT_s$ e
$0$ nos demais instantes de amostragem: as caudas se cancelam entre si.
\begin{center}
\begin{tikzpicture}
\begin{groupplot}[group style={group name=g, group size=3 by 1, horizontal sep=1.3cm}, ftdemot, height=0.18\textheight]
\nextgroupplot[title={(a) Sinal amostrado\\[1pt]{\tiny\ftstrut $x_\delta(t)$}}]
  \addplot[gray!55, dashed, domain=0:\fttmax, samples=300] {cos(deg(2*pi*\ftfzero*x))};
  \addplot[ycomb, ftazul, thick] coordinates {\ftamostras};
  \addplot[ftimppos, ftazul] coordinates {\ftamospos};
  \addplot[ftimpneg, ftazul] coordinates {\ftamosneg};
\nextgroupplot[title={(b) Núcleos de interpolação\\[1pt]{\tiny\ftstrut $x_f(nT_s)\,\mathrm{sinc}(f_s(t-nT_s))$}}, clip=true]
  \idnucleos
\nextgroupplot[title={(c) Soma dos núcleos\\[1pt]{\tiny\ftstrut $x_r(t) = x_f(t)$}}]
  \draw[black!45, densely dashed, thin] (axis cs:0,1) -- (axis cs:\fttmax,1);
  \addplot[gray!55, dashed, domain=0:\fttmax, samples=300] {cos(deg(2*pi*\ftfzero*x))};
  \addplot[ftverde, thick, densely dashed, domain=0:\fttmax, samples=300, visible on=<3->]
    {cos(deg(2*pi*\ftfzero*x))};
\end{groupplot}
\end{tikzpicture}\par\vspace{2pt}
{\scriptsize
\tikz[baseline=-0.6ex]\draw[ftazul, thick] (0,0) -- (0.4,0);\; amostras $x_f(nT_s)$
\qquad
\tikz[baseline=-0.6ex]\draw[ftroxo!55] (0,0) -- (0.4,0);\; núcleos $\mathrm{sinc}$
\qquad
\tikz[baseline=-0.6ex]\draw[ftverde, thick, densely dashed] (0,0) -- (0.4,0);\; sinal reconstruído
\qquad
\tikz[baseline=-0.6ex]\draw[gray!55, dashed] (0,0) -- (0.4,0);\; $x_f(t)$}
\end{center}
\pause\pause
\end{frame}

% ----------------------------------------------------------------------------
%  D. VISÃO GERAL DO PROCESSO + GANCHO
% ----------------------------------------------------------------------------
\begin{frame}{Método Ideal}{Visão Geral no Domínio do Tempo}
\centering
\begin{tikzpicture}
\begin{groupplot}[
  group style={group size=3 by 1, horizontal sep=1.2cm},
  ftaxis, width=0.26\textwidth
]

% (a) Sinal original
\nextgroupplot[title={(a) Sinal original\\[1pt]{\tiny\ftstrut $x_f(t)=\cos(2\pi\cdot 2t)$}}]
  \addplot[ftazul, thick, domain=0:\fttmax, samples=400]
    {cos(deg(2*pi*\ftfzero*x))};

% (b) Sinal amostrado (impulsos ponderados)
\nextgroupplot[title={(b) Sinal amostrado\\[1pt]{\tiny\ftstrut $x_\delta(t)$}}]
  \addplot[gray!55, dashed, domain=0:\fttmax, samples=400]
    {cos(deg(2*pi*\ftfzero*x))};
  \addplot[ycomb, ftazul, thick] coordinates {\ftamostras};
  \addplot[ftimppos, ftazul] coordinates {\ftamospos};
  \addplot[ftimpneg, ftazul] coordinates {\ftamosneg};

% (c) Após o passa-baixas: sinal idêntico ao original
\nextgroupplot[title={(c) Após o passa-baixas\\[1pt]{\tiny\ftstrut $x_r(t) = x_f(t)$}},
  clip mode=individual]
  \draw[black!45, densely dashed, thin] (axis cs:0,1) -- (axis cs:\fttmax,1);
  \addplot[gray!55, dashed, domain=0:\fttmax, samples=400]
    {cos(deg(2*pi*\ftfzero*x))};
  \addplot[ftverde, thick, domain=0:\fttmax, samples=400]
    {cos(deg(2*pi*\ftfzero*x))};
  \draw[gray, densely dotted] (axis cs:\ftpico,1) -- (axis cs:\ftpico,1.14);
  \node[font=\tiny, anchor=south west, inner sep=1pt] at (axis cs:\ftpico,1.12) {sem atraso};

\end{groupplot}
\end{tikzpicture}

\vspace{0.4em}
\mbox{\scriptsize
\tikz[baseline=-0.6ex]\draw[gray!55, dashed] (0,0) -- (0.45,0); $x_f(t)$\hspace{0.7em}
\tikz[baseline=-0.6ex]\draw[ftazul, thick] (0,0) -- (0.45,0); impulsos de área $x_f(nT_s)$\hspace{0.7em}
\tikz[baseline=-0.6ex]\draw[ftverde, thick] (0,0) -- (0.45,0); após o passa-baixas}

\vspace{0.4em}
{\small Amplitude preservada e nenhum atraso: não há equalizador a projetar.}
\end{frame}

\begin{frame}{Método Ideal}{Visão Geral no Domínio da Frequência}
\centering
\begin{tikzpicture}
\begin{groupplot}[
  group style={group size=3 by 1, horizontal sep=1.2cm},
  ftfreqaxis, width=0.26\textwidth
]

% (a) Espectro do sinal original
\nextgroupplot[title={(a) Espectro original\\[1pt]{\tiny\ftstrut $X_f(f)=\frac{1}{2}[\delta(f-2)+\delta(f+2)]$}}]
  \addplot[ftimpulso, ftazul] coordinates {(-\ftfzero,0.5) (\ftfzero,0.5)};

% (b) Espectro amostrado: réplicas em n f_s, todas com a mesma amplitude
\nextgroupplot[title={(b) Sinal amostrado\\[1pt]{\tiny\ftstrut $X_\delta(f)$}}]
  \draw[black!55, densely dashed] (axis cs:-\ftfmax,0.5) -- (axis cs:\ftfmax,0.5);
  \addplot[ftimpulso, ftlaranja!90!black] coordinates {\idrepmeio};
  \addplot[ftimpulso, ftazul] coordinates {(-\ftfzero,0.5) (\ftfzero,0.5)};

% (c) Após o passa-baixas: só a banda base, sem atenuação
\nextgroupplot[title={(c) Após o passa-baixas\\[1pt]{\tiny\ftstrut $X_r(f) = X_f(f)$}}]
  \draw[black!45, densely dashed, thin] (axis cs:-\ftfmax,0.5) -- (axis cs:\ftfmax,0.5);
  \draw[ftroxo, densely dashed, fill=ftroxo!8] (axis cs:-\ftfc,0) rectangle (axis cs:\ftfc,0.6);
  \node[font=\tiny, anchor=south, inner sep=1pt, yshift=4pt, ftroxo] at (axis cs:0,0.6) {$H_{PB}(f)$};
  \addplot[ftimpulso, ftverde] coordinates {(-\ftfzero,0.5) (\ftfzero,0.5)};

\end{groupplot}
\end{tikzpicture}

\vspace{0.4em}
\mbox{\scriptsize
\tikz[baseline=-0.6ex]\draw[black!55, densely dashed] (0,0) -- (0.45,0); envelope constante $\frac{1}{2}$\hspace{0.7em}
\tikz[baseline=-0.6ex]\draw[ftazul, thick] (0,0) -- (0.45,0); banda base\hspace{0.7em}
\tikz[baseline=-0.6ex]\draw[ftlaranja!90!black, thick] (0,0) -- (0.45,0); réplicas em $nf_s \pm 2$\hspace{0.7em}
\tikz[baseline=-0.6ex]\draw[ftroxo, densely dashed, fill=ftroxo!8] (0,-0.08) rectangle (0.5,0.12);\hspace{3pt}passa-baixas}

\vspace{0.4em}
{\small Envelope constante: todas as réplicas com a mesma amplitude, portanto nenhuma
distorção na banda base.}
\end{frame}

\section{Amostragem Natural}

\begin{frame}{Amostragem Natural}{Contexto}
	
	A amostragem ideal utiliza um trem de impulsos de Dirac $\delta(t)$, que é uma abstração matemática puramente teórica e fisicamente inexequível, pois exige pulsos com largura infinitesimal ($\tau \to 0$) e amplitude infinita para reter energia finita. 
	
	A amostragem natural surge para tornar o processo fisicamente realizável: em vez de impulsos instantâneos, utiliza-se um trem de pulsos retangulares de largura finita $\tau$ e período $T_s$, gerados por chaves eletrônicas reais (como transistores operando em comutação) que se fecham por um intervalo $\tau$ a cada ciclo. Ela aproxima o sistema de um cenário real de chaveamento analógico.
	
\end{frame}

% Definição do Pulso p0(t)
\begin{frame}{Amostragem Natural}{Definição do Pulso Fundamental $p_0(t)$}
	
	O processo de amostragem natural utiliza pulsos retangulares de largura finita $\tau$. Definimos inicialmente o pulso unitário $p_0(t)$ centrado na origem:
	
	\begin{equation*}
		p_0(t) = 
		\begin{cases}
			1, & -\frac{\tau}{2} \le t \le \frac{\tau}{2} \\[6pt]
			0, & \text{caso contrário}
		\end{cases}
	\end{equation*}
\end{frame}

\begin{frame}{Amostragem Natural}{Definição do Pulso Fundamental $p_0(t)$}
	\vfill
	\centering
	\begin{tikzpicture}
		\begin{axis}[
			pamaxis, 
			ylabel={$p_0(t)$}, 
			xmin=-0.6, xmax=0.6, 
			ymin=-0.2, ymax=1.4,
			xlabel={$t$}
			]
			% Pulso p0(t)
			\draw[sinalcor, line width=1.2pt, fill=sinalcor!20] 
			(axis cs:-0.2, 0) -- (axis cs:-0.2, 1) -- (axis cs:0.2, 1) -- (axis cs:0.2, 0) -- cycle;
			
			% Marcações do eixo t (-tau/2 e tau/2)
			\node[below, font=\small] at (axis cs:-0.2, 0) {$-\frac{\tau}{2}$};
			\node[below, font=\small] at (axis cs:0.2, 0) {$\frac{\tau}{2}$};
			\node[left, font=\small] at (axis cs:0, 1.1) {$1$};
		\end{axis}
	\end{tikzpicture}
\end{frame}

% Definição do Trem de Pulsos p(t)
\begin{frame}{Amostragem Natural}{Trem Periódico de Pulsos $p(t)$}
	
	O sinal de amostragem $p(t)$ é construído pela replicação periódica do pulso $p_0(t)$ a cada período de amostragem $T_s$:
	
	\begin{equation*}
		p(t) = \sum_{k=-\infty}^{\infty} p_0(t - k T_s)
	\end{equation*}
	
\end{frame}

\begin{frame}{Amostragem Natural}{Trem Periódico de Pulsos $p(t)$}
	\vfill
	\centering
	\begin{tikzpicture}
		\begin{axis}[pamaxis, ylabel={$p(t)$}, ymin=-0.2, ymax=1.4]
			\addplot[draw=none] {0};
			\pgfplotsextra{
				\foreach \n in {0,...,\Nam}{
					\pgfmathsetmacro\xa{\n*\Ts - \Tau/2}
					\pgfmathsetmacro\xb{\n*\Ts + \Tau/2}
					\draw[sinalcor, line width=1pt, fill=sinalcor!20] 
					(axis cs:\xa,0) rectangle (axis cs:\xb,1);
				}
			}
			\marcaTs
		\end{axis}
	\end{tikzpicture}
	\par\vspace{0.3em}
	{\small Trem de pulsos periódico com período $T_s$ e largura de pulso $\tau$.}
\end{frame}

% Conexão e Série de Fourier de p(t)
\begin{frame}{Amostragem Natural}{Expansão em Série de Fourier de $p(t)$}
	Como $p(t)$ é estritamente periódico com período $T_s$ ($f_s = 1/T_s$), podemos representá-lo via Série Exponencial de Fourier:
	\begin{equation*}
		p(t) = \sum_{k=-\infty}^{\infty} P[k] e^{j 2\pi k f_s t}
	\end{equation*}
	
	O cálculo dos coeficientes $P[k]$ é realizado no intervalo de um período $[-\frac{T_s}{2}, \frac{T_s}{2}]$:
	\begin{align*}
		P[k] &= \frac{1}{T_s} \int_{-T_s/2}^{T_s/2} p_0(t) e^{-j 2\pi k f_s t} \, dt = \frac{1}{T_s} \int_{-\tau/2}^{\tau/2} (1) e^{-j 2\pi k f_s t} \, dt \\[6pt]
		&= \frac{1}{T_s} \left[ \frac{e^{-j 2\pi k f_s t}}{-j 2\pi k f_s} \right]_{-\tau/2}^{\tau/2} = \frac{1}{T_s} \left( \frac{e^{j \pi k f_s \tau} - e^{-j \pi k f_s \tau}}{j 2\pi k f_s} \right)
	\end{align*}
\end{frame}

\begin{frame}{Amostragem Natural}{Coeficientes $P[k]$ da Série de Fourier}
	Aplicando a identidade de Euler $\sin(\theta) = \frac{e^{j\theta} - e^{-j\theta}}{2j}$:
	\begin{equation*}
		P[k] = \frac{\sin(\pi k f_s \tau)}{\pi k T_s f_s}
	\end{equation*}
	
	Como $T_s f_s = 1$, multiplicamos e dividimos por $\tau$ para obter a função $\text{sinc}$:
	\begin{equation*}
		P[k] = \frac{\tau}{T_s} \cdot \frac{\sin(\pi k f_s \tau)}{\pi k f_s \tau} \implies P[k] = \frac{\tau}{T_s} \text{sinc}(k f_s \tau)
	\end{equation*}
	
	\begin{block}{Sinal de Amostragem no Tempo}
		\begin{equation*}
			p(t) = \sum_{k=-\infty}^{\infty} \left[ \frac{\tau}{T_s} \text{sinc}(k f_s \tau) \right] e^{j 2\pi k f_s t}
		\end{equation*}
	\end{block}
\end{frame}

% Construção Analítica de s(t) no Tempo
\begin{frame}{Amostragem Natural}{Modelagem Analítica no Tempo}
	O sinal amostrado $s(t)$ é gerado pelo produto direto entre $x(t)$ e $p(t)$:
	
	\begin{equation*}
		s(t) = x(t) \cdot p(t)
	\end{equation*}
	
	Substituindo a expansão em Série de Fourier de $p(t)$:
	
	\begin{equation*}
		s(t) = x(t) \cdot \sum_{k=-\infty}^{\infty} P[k] e^{j 2\pi k f_s t}
	\end{equation*}
	
	
\end{frame}

\begin{frame}{Amostragem Natural}{Modelagem Analítica no Tempo}
	Pela propriedade distributiva, inserimos $x(t)$ no somatório:
	
	\begin{equation*}
		s(t) = \sum_{k=-\infty}^{\infty} P[k] x(t) e^{j 2\pi k f_s t}
	\end{equation*}
	
	
	Substituindo a expressão analítica de $P[k]$:
	
	\begin{equation*}
		s(t) = \sum_{k=-\infty}^{\infty} \left[ \frac{\tau}{T_s} \text{sinc}(k f_s \tau) \right] x(t) e^{j 2\pi k f_s t}
	\end{equation*}
\end{frame}

% Amostragem Natural: sobreposição x(t) + p(t), e produto s(t) = x(t)·p(t)
\begin{frame}{Amostragem Natural}{Sobreposição e Produto}
	\centering
	\begin{minipage}[c][0.72\textheight][c]{0.95\textwidth}
		\centering
		
		% Parâmetros locais (Ts = 0.2 s, Tau = 0.05 s, 5 pulsos)
		\def\Ts{0.2}
		\def\Tau{0.05}
		\def\Nam{5}
		
		\only<1>{
			\begin{tikzpicture}
				\begin{axis}[
					pamaxis,
					ylabel={$x(t)$ e $p(t)$},
					ylabel style={at={(axis cs:0,1.45)}, anchor=south}, % Centraliza sobre o eixo x=0
					trig format=rad,
					ymin=-0.2, ymax=1.45
					]
					% Trem de pulsos p(t)
					\pgfplotsextra{%
						\foreach \n in {0,...,\Nam}{%
							\pgfmathsetmacro\xa{\n*\Ts - \Tau/2}%
							\pgfmathsetmacro\xb{\n*\Ts + \Tau/2}%
							\draw[gray!60, line width=.7pt, fill=gray!15]
							(axis cs:\xa,0) rectangle (axis cs:\xb,1);
						}%
					}
					% x(t) = cos(2*pi*2*t)
					\addplot[sinalcor, line width=1.4pt, smooth, domain=0:1, samples=300]
					{cos(4*pi*x)};
				\end{axis}
			\end{tikzpicture}\par
			\small Sinal analógico $x(t)=\cos(2\pi\cdot 2t)$ sobreposto ao trem de pulsos $p(t)$ (largura $\tau$, período $T_s$).}
		
		% produto s(t)
		\only<2>{%
			\begin{tikzpicture}
				\begin{axis}[
					pamaxis,
					ylabel={$s(t) = x(t) \cdot p(t)$},
					ylabel style={at={(axis cs:0,1.45)}, anchor=south}, % Centraliza sobre o eixo x=0
					trig format=rad,
					ymin=-0.2, ymax=1.45
					]
					% Envoltória x(t) como referência tracejada
					\addplot[gray!60, dashed, line width=.7pt, smooth, domain=0:1, samples=300]
					{cos(4*pi*x)};
					% Pulsos de amostragem natural: topo SEGUE x(t)
					\pgfplotsextra{%
						\foreach \n in {0,...,\Nam}{%
							\pgfmathsetmacro\xa{\n*\Ts - \Tau/2}%
							\pgfmathsetmacro\xb{\n*\Ts + \Tau/2}%
							\draw[sinalcor, line width=.9pt, fill=sinalcor!25]
							(axis cs:\xa,0) --
							(axis cs:\xa,{cos(4*pi*\xa)}) --
							plot[domain=\xa:\xb, samples=25, smooth]
							(axis cs:\x,{cos(4*pi*\x)}) --
							(axis cs:\xb,{cos(4*pi*\xb)}) --
							(axis cs:\xb,0) -- cycle;
						}%
					}
				\end{axis}
			\end{tikzpicture}\par
			\small Sinal amostrado naturalmente $s(t)=x(t)\cdot p(t)$. O topo de cada pulso \textbf{acompanha a forma de $x(t)$} — característica da amostragem natural.}
		
		% ---------------------------------------------------- (3) destaque: onde p(t)=0 vs p(t)=1
		\only<3>{%
			\begin{tikzpicture}
				\begin{axis}[
					pamaxis,
					ylabel={$s(t)$},
					ylabel style={at={(axis cs:0,1.45)}, anchor=south}, % Centraliza sobre o eixo x=0
					trig format=rad,
					ymin=-0.2, ymax=1.45
					]
					% x(t) como referência tracejada
					\addplot[gray!60, dashed, line width=.7pt, smooth, domain=0:1, samples=300]
					{cos(4*pi*x)};
					% Regiões onde p(t) = 0  ->  s(t) = 0 (faixas cinzas claras)
					\pgfplotsextra{%
						\foreach \n in {0,...,\Nam}{%
							\pgfmathsetmacro\xa{\n*\Ts + \Tau/2}%
							\pgfmathsetmacro\xb{(\n+1)*\Ts - \Tau/2}%
							\ifnum\n<\Nam
							\draw[black!10, fill=black!5]
							(axis cs:\xa,-0.15) rectangle (axis cs:\xb,1.35);
							\fi
						}%
					}
					% Pulsos de amostragem natural
					\pgfplotsextra{%
						\foreach \n in {0,...,\Nam}{%
							\pgfmathsetmacro\xa{\n*\Ts - \Tau/2}%
							\pgfmathsetmacro\xb{\n*\Ts + \Tau/2}%
							\draw[sinalcor, line width=.9pt, fill=sinalcor!25]
							(axis cs:\xa,0) --
							(axis cs:\xa,{cos(4*pi*\xa)}) --
							plot[domain=\xa:\xb, samples=25, smooth]
							(axis cs:\x,{cos(4*pi*\x)}) --
							(axis cs:\xb,{cos(4*pi*\xb)}) --
							(axis cs:\xb,0) -- cycle;
						}%
					}
					% Rótulos das duas regiões
					\node[font=\tiny, anchor=north] at (axis cs:0.30,-0.02)
					{$p(t)=0 \Rightarrow s(t)=0$};
					\node[font=\tiny, anchor=north, text=sinalcor] at (axis cs:0.05,-0.02)
					{$p(t)=1$};
				\end{axis}
			\end{tikzpicture}\par
			\small Onde $p(t)=0$, o produto é nulo (faixas cinzas). Onde $p(t)=1$ (dentro de cada pulso), o produto reproduz $x(t)$ — o topo \textbf{segue} o cosseno.}
		
	\end{minipage}
\end{frame}

% Domínio da Frequência: Convolução e Espectro Final
\begin{frame}{Amostragem Natural}{Análise no Domínio da Frequência}
	
	A multiplicação no tempo equivale à convolução no domínio da frequência:
	
	\begin{equation*}
		S(f) = \mathcal{F}\{s(t)\} = X(f) * P(f)
	\end{equation*}
	
	A Transformada de Fourier do trem de pulsos $p(t)$ é:
	
	\begin{equation*}
		P(f) = \sum_{k=-\infty}^{\infty} P[k] \delta(f - k f_s)
	\end{equation*}
	
\end{frame}

\begin{frame}{Amostragem Natural}{Análise no Domínio da Frequência}
	Calculando a convolução $X(f) * P(f)$:
	\begin{align*}
		S(f) &= X(f) * \sum_{k=-\infty}^{\infty} P[k] \delta(f - k f_s) \\[4pt]
		&= \sum_{k=-\infty}^{\infty} P[k] \left[ X(f) * \delta(f - k f_s) \right] \\[4pt]
		&= \sum_{k=-\infty}^{\infty} P[k] X(f - k f_s)
	\end{align*}
\end{frame}

% Frame 2: Espectro Final do Sinal Amostrado
\begin{frame}{Amostragem Natural}{Espectro Final do Sinal Amostrado}
	
	Substituindo o coeficiente $P[k] = \frac{\tau}{T_s} \text{sinc}(k f_s \tau)$, temos a expressão espectral final:
	
	\begin{equation*}
		S(f) = \sum_{k=-\infty}^{\infty} \frac{\tau}{T_s} \text{sinc}(k f_s \tau) X(f - k f_s)
	\end{equation*}
	
	\vfill
	\centering
	\begin{tikzpicture}
		\begin{axis}[
			pamaxis, 
			ylabel={$S(f)$}, 
			xlabel={$f$},
			xmin=-15, xmax=15, 
			ymin=-0.1, ymax=0.6,
			axis lines=middle,
			trig format=rad
			]
			% Envelope Sinc
			% Réplicas Espectrais centradas em k*fs
		\end{axis}
	\end{tikzpicture}
	\par\vspace{0.3em}
	{\small Réplicas de $X(f)$ repetidas a cada $f_s$, atenuadas pelo envelope $\frac{\tau}{T_s}\text{sinc}(k f_s \tau)$.}
\end{frame}
\begin{frame}{Método Topo-Plano: Amostragem}{Análise no Domínio do Tempo}

    Multiplicamos o sinal original $x(t)$ por um trem de impulsos deslocados, assim como na amostragem ideal.

\end{frame}

\begin{frame}{Método Topo-Plano: Amostragem}{Análise no Domínio do Tempo}

    A equação para a modulação é dada por:
    \begin{equation*}
        x_p(t) = \left[ x(t) \cdot \sum_{n=-\infty}^{\infty} \delta(t - nT_s) \right] * h(t)
    \end{equation*}

    Onde $h(t)$ é o pulso retangular de duração $\tau$:
    \begin{equation*}
        h(t) =
        \begin{cases}
            1, & 0 \leq t \leq \tau \\
            0, & \text{caso contrário.}
        \end{cases}
    \end{equation*}

\end{frame}

\begin{frame}{Método Topo-Plano: Amostragem}{Interpretação Gráfica}
\centering
\begin{minipage}[c][0.66\textheight][c]{0.92\textwidth}
\centering

% ------------------------------------------------------- (1) sinal original
\only<1>{%
\begin{tikzpicture}
    \begin{axis}[pamaxis, ylabel={$x(t)$}]
        \addplot[sinalcor, line width=1pt, smooth] {\sinal};
    \end{axis}
\end{tikzpicture}\par
Sinal analógico de banda base $x(t)$.}

% ------------------------------------------------------ (2) trem de impulsos
\only<2>{%
\begin{tikzpicture}
    \begin{axis}[pamaxis, ylabel={$\delta_{T_s}(t)$}]
        \addplot[draw=none] {0};
        \pgfplotsextra{%
            \foreach \n in {0,...,\Nam}{%
                \pgfmathsetmacro\xn{\n*\Ts}%
                \draw[sinalcor, line width=1pt, -{Stealth[length=5pt,width=4pt]}]
                    (axis cs:\xn,0) -- (axis cs:\xn,0.85);
                \draw[gray, line width=.4pt] (axis cs:\xn,-0.07) -- (axis cs:\xn,0.07);}}
        \marcaTs
    \end{axis}
\end{tikzpicture}\par
Trem de impulsos com período $T_s$.}

% ------------------------------------------- (3) produto (amostragem ideal)
\only<3>{%
\begin{tikzpicture}
    \begin{axis}[pamaxis, ylabel={$x_\delta(t)$}]
        \addplot[gray!50, dashed, line width=.7pt, smooth] {\sinal};
        \pgfplotsextra{%
            \foreach \n in {0,...,\Nam}{%
                \pgfmathsetmacro\xn{\n*\Ts}%
                \pgfmathsetmacro\yn{sin(deg(2*pi*0.75*\n*\Ts + 0.4))}%
                \draw[sinalcor, line width=1pt, -{Stealth[length=5pt,width=4pt]}]
                    (axis cs:\xn,0) -- (axis cs:\xn,\yn);}}
    \end{axis}
\end{tikzpicture}\par
Produto: impulsos ponderados por $x(nT_s)$.}

% ------------------------------------------------------ (4) pulso h(t)
\only<4>{%
\begin{tikzpicture}
    \begin{axis}[pamaxis, ylabel={$h(t)$}, ymin=-0.5, ymax=1.8,
                 xlabel style={at={(axis description cs:1,0.217)}, anchor=west}]
        \addplot[draw=none] {0};
        \pgfplotsextra{%
            \draw[sinalcor, line width=1pt, fill=sinalcor!18]
                (axis cs:0,0) rectangle (axis cs:\Tau,1);
            \draw[sinalcor, line width=1pt] (axis cs:\Tau,0) -- (axis cs:1.12,0);
            \node[below=1pt, font=\small] at (axis cs:\Tau,0) {$\tau$};
            \node[left, font=\small] at (axis cs:0,1) {$1$};}
    \end{axis}
\end{tikzpicture}\par
Pulso de retenção $h(t)$, de duração $\tau$.}

% --------------------------------------------------- (5) PAM de topo plano
\only<5>{%
\begin{tikzpicture}
    \begin{axis}[pamaxis, ylabel={$x_p(t)$}]
        \addplot[gray!50, dashed, line width=.7pt, smooth] {\sinal};
        \pgfplotsextra{%
            \foreach \n in {0,...,\Nam}{%
                \pgfmathsetmacro\xa{\n*\Ts}%
                \pgfmathsetmacro\xb{\n*\Ts+\Tau}%
                \pgfmathsetmacro\yn{sin(deg(2*pi*0.75*\n*\Ts + 0.4))}%
                \draw[sinalcor, line width=.9pt, fill=sinalcor!18]
                    (axis cs:\xa,0) rectangle (axis cs:\xb,\yn);}}
        \marcaTs
    \end{axis}
\end{tikzpicture}\par
Resultado da convolução: PAM com topo plano ($d = \tau/T_s = 1/3$).}

\end{minipage}
\end{frame}

% ============================================================================
%  Efeito de abertura e deslocamento: sinal original, amostrado,
%  após o passa-baixas e após o equalizador.
%  Parâmetros e macros (prefixo \ft...) definidos no main.tex.
% ============================================================================


% ============================================================================
%  Domínio da frequência: espectro original, amostrado,
%  após o passa-baixas e após o equalizador.
% ============================================================================


% ============================================================================
%  Demonstração no domínio do tempo (3 slides)
%  Azul: etapa 1 (multiplicação pelo trem de impulsos)
%  Laranja: etapa 2 (convolução com o pulso h)
% ============================================================================
\begin{frame}{Método Topo-Plano: Amostragem}{Demonstração no Domínio do Tempo}
\small
No domínio do tempo, o processo é representado como a \ftamos{multiplicação do sinal
original por um trem de impulsos}, seguida da \ftret{convolução desse resultado com um
pulso de largura $\tau$}, conforme a equação de modulação:
\begin{equation*}
  x_p(t) = \ftamos{\Big[\, x(t)\cdot p(t) \,\Big]}\;\ftret{*\; h(t)},
  \qquad
  p(t) = \sum_{n=-\infty}^{\infty} \delta(t - nT_s),
  \qquad
  h(t) =
  \begin{cases}
    1, & 0 \leq t \leq \tau \\
    0, & \text{caso contrário.}
  \end{cases}
\end{equation*}
\begin{center}
\begin{tikzpicture}
\begin{groupplot}[group style={group name=g, group size=3 by 1, horizontal sep=1.3cm}, ftdemot]
\nextgroupplot[title={(a) Sinal original\\[1pt]{\tiny\ftstrut $x(t)$}}]
  \addplot[ftazul!40!black, thick, domain=0:\fttmax, samples=300] {cos(deg(2*pi*\ftfzero*x))};
\nextgroupplot[title={(b) Trem de impulsos\\[1pt]{\tiny\ftstrut $p(t)$}}]
  \addplot[ycomb, ftazul, thick] coordinates {\ftimpt};
  \addplot[ftimppos, ftazul] coordinates {\ftimpt};
  \draw[gray, <->, >=stealth] (axis cs:\ftTs,1.18) -- node[above, font=\tiny, inner sep=1pt]{$T_s$} (axis cs:2*\ftTs,1.18);
\nextgroupplot[title={(c) Pulso retangular\\[1pt]{\tiny\ftstrut $h(t)$}}, xmin=0, xmax=0.5,
  xtick={0,\fttau,\ftTs}, xticklabels={$0$,$\tau$,$T_s$}]
  \draw[ftlaranja!90!black, thick, fill=ftlaranja!22] (axis cs:0,0) rectangle (axis cs:\fttau,1);
\end{groupplot}
\end{tikzpicture}
\end{center}
\end{frame}

\begin{frame}{Método Topo-Plano: Amostragem}{Demonstração no Domínio do Tempo}
\small
Realizando a multiplicação, obtemos o trem de impulsos ponderado pelo sinal original,
como na amostragem ideal:
\begin{equation*}
  x_p(t) = \ftamos{\left[\, \sum_{n=-\infty}^{\infty} x(nT_s)\,\delta(t - nT_s) \right]}\;\ftret{*\; h(t)}
\end{equation*}
\begin{center}
\begin{tikzpicture}
\begin{groupplot}[group style={group name=g, group size=3 by 1, horizontal sep=1.3cm}, ftdemot]
\nextgroupplot[title={(a) Sinal original\\[1pt]{\tiny\ftstrut $x(t)$}}]
  \addplot[ftazul!40!black, thick, domain=0:\fttmax, samples=300] {cos(deg(2*pi*\ftfzero*x))};
\nextgroupplot[title={(b) Trem de impulsos\\[1pt]{\tiny\ftstrut $p(t)$}}]
  \addplot[ycomb, ftazul, thick] coordinates {\ftimpt};
  \addplot[ftimppos, ftazul] coordinates {\ftimpt};
  \draw[gray, <->, >=stealth] (axis cs:\ftTs,1.18) -- node[above, font=\tiny, inner sep=1pt]{$T_s$} (axis cs:2*\ftTs,1.18);
\nextgroupplot[title={(c) Amostragem ideal\\[1pt]{\tiny\ftstrut $x(t)\,p(t)$}}]
  \addplot[gray!55, dashed, domain=0:\fttmax, samples=300] {cos(deg(2*pi*\ftfzero*x))};
  \addplot[ycomb, ftazul, thick] coordinates {\ftamostras};
  \addplot[ftimppos, ftazul] coordinates {\ftamospos};
  \addplot[ftimpneg, ftazul] coordinates {\ftamosneg};
\end{groupplot}
\ftop{1}{2}{\times}
\ftop{2}{3}{=}
\end{tikzpicture}
\end{center}
\end{frame}

\begin{frame}{Método Topo-Plano: Amostragem}{Demonstração no Domínio do Tempo}
\small
Calculando a convolução com o pulso, obtemos o sinal amostrado final: cada impulso é
substituído por um pulso de largura $\tau$, pois $\delta(t - nT_s) * h(t) = h(t - nT_s)$:
\begin{equation*}
  \boxed{\;
  x_p(t) = \sum_{n=-\infty}^{\infty} x(nT_s)\, h(t - nT_s)
  \;}
\end{equation*}
\begin{center}
\begin{tikzpicture}
\begin{groupplot}[group style={group name=g, group size=3 by 1, horizontal sep=1.3cm}, ftdemot]
\nextgroupplot[title={(a) Amostragem ideal\\[1pt]{\tiny\ftstrut $x(t)\,p(t)$}}]
  \addplot[gray!55, dashed, domain=0:\fttmax, samples=300] {cos(deg(2*pi*\ftfzero*x))};
  \addplot[ycomb, ftazul, thick] coordinates {\ftamostras};
  \addplot[ftimppos, ftazul] coordinates {\ftamospos};
  \addplot[ftimpneg, ftazul] coordinates {\ftamosneg};
\nextgroupplot[title={(b) Pulso retangular\\[1pt]{\tiny\ftstrut $h(t)$}}, xmin=0, xmax=0.5,
  xtick={0,\fttau,\ftTs}, xticklabels={$0$,$\tau$,$T_s$}]
  \draw[ftlaranja!90!black, thick, fill=ftlaranja!22] (axis cs:0,0) rectangle (axis cs:\fttau,1);
\nextgroupplot[title={(c) Sinal amostrado\\[1pt]{\tiny\ftstrut $x_p(t)$}}]
  \addplot[gray!55, dashed, domain=0:\fttmax, samples=300] {cos(deg(2*pi*\ftfzero*x))};
  \addplot[ftlaranja!90!black, semithick, fill=ftlaranja!22] coordinates {\ftpulsos};
\end{groupplot}
\ftop{1}{2}{*}
\ftop{2}{3}{=}
\end{tikzpicture}
\end{center}
\end{frame}

% ============================================================================
%  Demonstração no domínio da frequência (3 slides)
% ============================================================================
\begin{frame}{Método Topo-Plano: Amostragem}{Demonstração no Domínio da Frequência}
\small
Para obter o espectro do sinal amostrado, aplicamos a transformada de Fourier à equação
de modulação. O produto no tempo vira convolução e a convolução vira produto:
\begin{equation*}
  x_p(t) = \ftamos{\Big[\, x(t)\cdot p(t) \,\Big]}\;\ftret{*\; h(t)}
  \quad\overset{\mathcal{F}}{\longleftrightarrow}\quad
  X_p(f) = \ftamos{\Big[\, X(f) * P(f) \,\Big]}\;\ftret{\cdot\; H(f)}
\end{equation*}
\vspace{-0.6em}
\begin{equation*}
  \ftamos{P(f) = \frac{1}{T_s}\sum_{n=-\infty}^{\infty} \delta(f - nf_s)},
  \qquad
  \ftret{H(f) = \tau\, \mathrm{sinc}(f\tau)\, e^{-j\pi f\tau}},
  \qquad
  \mathrm{sinc}(x) = \frac{\operatorname{sen}(\pi x)}{\pi x}
\end{equation*}
\begin{center}
\begin{tikzpicture}
\begin{groupplot}[group style={group name=g, group size=3 by 1, horizontal sep=1.3cm}, ftdemof]
\nextgroupplot[title={(a) Espectro original\\[1pt]{\tiny\ftstrut $X(f)$}}, yticklabels={$\frac{1}{2}$}]
  \addplot[ftimpulso, ftazul!40!black] coordinates {(-\ftfzero,1) (\ftfzero,1)};
\nextgroupplot[title={(b) Trem de impulsos\\[1pt]{\tiny\ftstrut $P(f)$}}, yticklabels={$\frac{1}{T_s}$}]
  \addplot[ftimpulso, ftazul] coordinates {\ftimpf};
\nextgroupplot[title={(c) Espectro do pulso\\[1pt]{\tiny\ftstrut $|H(f)|$}}, yticklabels={$\tau$}]
  \addplot[ftlaranja!90!black, thick, domain=-\ftfmax:\ftfmax, samples=400]
    {abs(sin(deg(pi*x*\fttau))/(pi*x*\fttau))};
  \node[font=\tiny, anchor=south west, inner sep=1pt] at (axis cs:\ftnulo,0.02) {$\frac{1}{\tau}$};
\end{groupplot}
\end{tikzpicture}
\end{center}
\end{frame}

\begin{frame}{Método Topo-Plano: Amostragem}{Demonstração no Domínio da Frequência}
\small
Ao calcular a convolução entre $X(f)$ e $P(f)$, cada impulso de $P(f)$ cria uma cópia
de $X(f)$, como na amostragem ideal:
\begin{equation*}
  \ftamos{X(f) * P(f) = \frac{1}{T_s} \sum_{n=-\infty}^{\infty} X(f - nf_s)}
\end{equation*}
\begin{center}
\begin{tikzpicture}
\begin{groupplot}[group style={group name=g, group size=3 by 1, horizontal sep=1.3cm}, ftdemof]
\nextgroupplot[title={(a) Espectro original\\[1pt]{\tiny\ftstrut $X(f)$}}, yticklabels={$\frac{1}{2}$}]
  \addplot[ftimpulso, ftazul!40!black] coordinates {(-\ftfzero,1) (\ftfzero,1)};
\nextgroupplot[title={(b) Trem de impulsos\\[1pt]{\tiny\ftstrut $P(f)$}}, yticklabels={$\frac{1}{T_s}$}]
  \addplot[ftimpulso, ftazul] coordinates {\ftimpf};
\nextgroupplot[title={(c) Amostragem ideal\\[1pt]{\tiny\ftstrut $X(f)*P(f)$}}, yticklabels={$\frac{1}{2T_s}$}]
  \addplot[ftimpulso, ftazul] coordinates {\ftrepideal};
\end{groupplot}
\ftop{1}{2}{*}
\ftop{2}{3}{=}
\end{tikzpicture}
\end{center}
\end{frame}

\begin{frame}{Método Topo-Plano: Amostragem}{Demonstração no Domínio da Frequência}
\small
Calculando, agora, a multiplicação entre o resultado da convolução e o espectro do pulso,
obtemos:
\begin{equation*}
  \boxed{\;
  X_p(f) = \frac{\tau}{T_s}\;
  \ftret{e^{-j\pi f\tau}\, \mathrm{sinc}(f\tau)}
  \;\ftamos{\sum_{n=-\infty}^{\infty} X(f - nf_s)}
  \;}
\end{equation*}
\begin{center}
\begin{tikzpicture}
\begin{groupplot}[group style={group name=g, group size=3 by 1, horizontal sep=1.3cm}, ftdemof]
\nextgroupplot[title={(a) Amostragem ideal\\[1pt]{\tiny\ftstrut $X(f)*P(f)$}}, yticklabels={$\frac{1}{2T_s}$}]
  \addplot[ftimpulso, ftazul] coordinates {\ftrepideal};
\nextgroupplot[title={(b) Espectro do pulso\\[1pt]{\tiny\ftstrut $|H(f)|$}}, yticklabels={$\tau$}]
  \addplot[ftlaranja!90!black, thick, domain=-\ftfmax:\ftfmax, samples=400]
    {abs(sin(deg(pi*x*\fttau))/(pi*x*\fttau))};
  \node[font=\tiny, anchor=south west, inner sep=1pt] at (axis cs:\ftnulo,0.02) {$\frac{1}{\tau}$};
\nextgroupplot[title={(c) Sinal amostrado\\[1pt]{\tiny\ftstrut $|X_p(f)|$}}, yticklabels={$\frac{\tau}{2T_s}$}]
  \addplot[black!55, densely dashed, domain=-\ftfmax:\ftfmax, samples=400]
    {abs(sin(deg(pi*x*\fttau))/(pi*x*\fttau))};
  \addplot[ftimpulso, ftlaranja!90!black] coordinates {\ftrepsinc};
\end{groupplot}
\ftop{1}{2}{\times}
\ftop{2}{3}{=}
\end{tikzpicture}
\end{center}
\end{frame}

% ============================================================================
%  Efeito de abertura (dinâmico: tau/T_s muda a cada clique)
% ============================================================================
\begin{frame}{Método Topo-Plano: Efeito de Abertura}{Representação e Demonstração}
\small
Após o passa-baixas ideal (ganho $T_s$, corte $f_s/2$), resta apenas a banda base,
multiplicada pelo espectro do pulso:
\begin{equation*}
  X_r'(f) = T_s\,X_p(f) = \tau\;\ftret{\mathrm{sinc}(f\tau)}\;e^{-j\pi f\tau}\;X(f),
  \qquad |f| < \frac{f_s}{2}
\end{equation*}
O fator $\ftret{\mathrm{sinc}(f\tau)}$ atenua cada componente conforme sua frequência.
Para $x(t)=\cos(2\pi\cdot 2t)$, a amplitude cai para $A = \mathrm{sinc}(2\tau)$.
\begin{center}
\foreach \d/\k in {0.25/1,0.5/2,0.75/3,1/4}{%
\only<\k>{%
\pgfmathsetmacro\tk{\d/\ftfs}%
\pgfmathsetmacro\Ak{sin(deg(pi*\ftfzero*\tk))/(pi*\ftfzero*\tk)}%
\pgfmathsetmacro\Bk{sin(deg(pi*\d/2))/(pi*\d/2)}%
\pgfmathsetmacro\dBk{20*log10(\Bk)}%
\begin{tikzpicture}
\begin{groupplot}[group style={group name=g, group size=3 by 1, horizontal sep=1.3cm}, ftdemot, height=0.18\textheight]
\nextgroupplot[title={(a) Espectro do pulso\\[1pt]{\tiny\ftstrut $|H(f)|/\tau$}},
  xmin=-\ftfmax, xmax=\ftfmax, ymin=0, ymax=1.3, xtick={-10,-5,0,5,10}, ytick={0,0.5,1},
  extra y ticks={}, xlabel={$f$ (Hz)}]
  \fill[ftroxo!8] (axis cs:-\ftfc,0) rectangle (axis cs:\ftfc,1.3);
  \addplot[ftlaranja!90!black, thick, domain=-\ftfmax:\ftfmax, samples=400]
    {abs(sin(deg(pi*x*\tk))/(pi*x*\tk))};
  \addplot[only marks, mark=*, mark size=1.5pt, ftazul] coordinates {(-\ftfzero,\Ak) (\ftfzero,\Ak)};
\nextgroupplot[title={(b) Atenuação\\[1pt]{\tiny\ftstrut em função de $\tau/T_s$}},
  xmin=0, xmax=1, ymin=0.6, ymax=1.05, xtick={0,0.25,0.5,0.75,1}, ytick={0.6,0.7,0.8,0.9,1},
  extra y ticks={}, xlabel={$\tau/T_s$}]
  \addplot[ftazul, thick, domain=0.001:1, samples=60] {sin(deg(pi*\ftfzero*x/\ftfs))/(pi*\ftfzero*x/\ftfs)};
  \addplot[ftlaranja!90!black, thick, densely dashed, domain=0.001:1, samples=60] {sin(deg(pi*x/2))/(pi*x/2)};
  \addplot[only marks, mark=*, mark size=1.6pt, ftazul] coordinates {(\d,\Ak)};
  \addplot[only marks, mark=*, mark size=1.6pt, ftlaranja!90!black] coordinates {(\d,\Bk)};
\nextgroupplot[title={(c) Amplitude no tempo\\[1pt]{\tiny\ftstrut sem o atraso, para isolar $A$}}]
  \draw[black!45, densely dashed, thin] (axis cs:0,1) -- (axis cs:\fttmax,1);
  \addplot[gray!55, dashed, domain=0:\fttmax, samples=300] {cos(deg(2*pi*\ftfzero*x))};
  \addplot[ftlaranja!90!black, thick, domain=0:\fttmax, samples=300] {\Ak*cos(deg(2*pi*\ftfzero*x))};
\end{groupplot}
\end{tikzpicture}\par\vspace{2pt}
{\scriptsize
$\tau/T_s = \pgfmathprintnumber[fixed,precision=2,use comma]{\d}$:\quad
\tikz[baseline=-0.6ex]\fill[ftazul] (0,0) circle (1.6pt);\;
$A = \mathrm{sinc}(f_0\tau) = \pgfmathprintnumber[fixed,precision=3,use comma]{\Ak}$
\qquad
\tikz[baseline=-0.6ex]\draw[ftlaranja!90!black, thick, densely dashed] (0,0) -- (0.4,0);\;
queda na borda da banda: $\mathrm{sinc}\!\left(\frac{\tau}{2T_s}\right) = \pgfmathprintnumber[fixed,precision=3,use comma]{\Bk}$
($\pgfmathprintnumber[fixed,precision=2,use comma]{\dBk}$~dB)}
}}
\end{center}
\end{frame}

% ============================================================================
%  Deslocamento (dinâmico: tau/T_s muda a cada clique)
% ============================================================================
\begin{frame}{Método Topo-Plano: Deslocamento}{Representação e Demonstração}
\small
O pulso causal é o pulso centrado atrasado de $\tau/2$. Pela propriedade do deslocamento
no tempo, $g(t-t_0) \leftrightarrow G(f)\,e^{-j2\pi f t_0}$, com $t_0 = \tau/2$:
\begin{equation*}
  h(t) = h_c\!\left(t - \frac{\tau}{2}\right)
  \;\overset{\mathcal{F}}{\longleftrightarrow}\;
  H(f) = \tau\,\mathrm{sinc}(f\tau)\;\ftret{e^{-j\pi f\tau}},
  \qquad
  \angle H(f) = -\pi f\tau
  \;\Rightarrow\;
  t_g = -\frac{1}{2\pi}\frac{d\angle H}{df} = \frac{\tau}{2}
\end{equation*}
A fase é linear: todas as frequências sofrem o mesmo atraso $\tau/2$, sem distorcer a forma do sinal.
\begin{center}
\foreach \d/\k in {0.25/1,0.5/2,0.75/3,1/4}{%
\only<\k>{%
\pgfmathsetmacro\tk{\d/\ftfs}%
\pgfmathsetmacro\atk{\d*500/\ftfs}%
\pgfmathsetmacro\pk{\ftpico+\tk/2}%
\begin{tikzpicture}
\begin{groupplot}[group style={group name=g, group size=3 by 1, horizontal sep=1.3cm}, ftdemot, height=0.18\textheight]
\nextgroupplot[title={(a) Pulso centrado e causal\\[1pt]{\tiny\ftstrut $h_c(t)$ e $h(t)$}},
  xmin=-0.15, xmax=0.28, xtick={-0.1,0,0.1,0.2}]
  \draw[gray, densely dashed, thick] (axis cs:-\tk/2,0) rectangle (axis cs:\tk/2,1);
  \draw[ftlaranja!90!black, thick, fill=ftlaranja!22] (axis cs:0,0) rectangle (axis cs:\tk,1);
  \draw[black!80, ->, >=stealth] (axis cs:0,1.12) -- (axis cs:\tk/2,1.12);
  \node[font=\tiny, anchor=south west, inner sep=1pt] at (axis cs:\tk/2,1.1) {$\tau/2$};
\nextgroupplot[title={(b) Fase do pulso\\[1pt]{\tiny\ftstrut $\angle H(f) = -\pi f\tau$}},
  xmin=-\ftfc, xmax=\ftfc, ymin=-1.8, ymax=1.8,
  xtick={-\ftfc,0,\ftfc}, xticklabels={$-\frac{f_s}{2}$,$0$,$\frac{f_s}{2}$},
  ytick={-1.5708,0,1.5708}, yticklabels={$-\frac{\pi}{2}$,$0$,$\frac{\pi}{2}$},
  extra y ticks={0}, xlabel={$f$ (Hz)}]
  \addplot[ftlaranja!90!black, thick, domain=-\ftfc:\ftfc, samples=2] {-pi*x*\tk};
\nextgroupplot[title={(c) Atraso no tempo\\[1pt]{\tiny\ftstrut sem a atenuação, para isolar $\tau/2$}}]
  \addplot[gray!55, dashed, domain=0:\fttmax, samples=300] {cos(deg(2*pi*\ftfzero*x))};
  \addplot[ftlaranja!90!black, thick, domain=0:\fttmax, samples=300] {cos(deg(2*pi*\ftfzero*(x-\tk/2)))};
  \draw[gray, densely dotted] (axis cs:\ftpico,1) -- (axis cs:\ftpico,1.18);
  \draw[ftlaranja!90!black, densely dotted] (axis cs:\pk,1) -- (axis cs:\pk,1.18);
\end{groupplot}
\end{tikzpicture}\par\vspace{2pt}
{\scriptsize
$\tau/T_s = \pgfmathprintnumber[fixed,precision=2,use comma]{\d}$:\quad
atraso $\tau/2 = \pgfmathprintnumber[fixed,precision=1,use comma]{\atk}$~ms
\qquad
\tikz[baseline=-0.6ex]\draw[gray, densely dashed, thick] (0,-0.08) rectangle (0.3,0.12);\; pulso centrado $h_c(t)$
\qquad
\tikz[baseline=-0.6ex]\draw[ftlaranja!90!black, thick, fill=ftlaranja!22] (0,-0.08) rectangle (0.3,0.12);\; pulso causal $h(t)$}
}}
\end{center}
\end{frame}

% ============================================================================
%  Reconstrução no domínio do tempo (revelado em 3 etapas)
% ============================================================================
\begin{frame}{Método Topo-Plano: Reconstrução}{Representação e Demonstração no Domínio do Tempo}
\small
Reconstruir é convoluir o sinal amostrado com $h_{PB}(t) = \mathrm{sinc}(f_s t)$, a resposta
ao impulso do passa-baixas ideal. Cada pulso vira um núcleo de interpolação
$g(t) = h(t) * h_{PB}(t)$ e, como $x(t)$ é limitado em banda, a soma é uma
\ftret{média móvel} do sinal original:
\begin{equation*}
  x_r'(t) = \sum_{n=-\infty}^{\infty} x(nT_s)\, g(t - nT_s)
  = x(t) * h(t) = \ftret{\int_{t-\tau}^{t} x(\lambda)\,d\lambda}
  \quad\xrightarrow{\text{equalizador}}\quad
  \boxed{\,x_r(t) = x\!\left(t - \frac{\tau}{2}\right)}
\end{equation*}
\begin{center}
\begin{tikzpicture}
\begin{groupplot}[group style={group name=g, group size=3 by 1, horizontal sep=1.3cm}, ftdemot, height=0.18\textheight]
\nextgroupplot[title={(a) Sinal amostrado\\[1pt]{\tiny\ftstrut $x_p(t)$}}]
  \addplot[gray!55, dashed, domain=0:\fttmax, samples=300] {cos(deg(2*pi*\ftfzero*x))};
  \addplot[ftlaranja!90!black, semithick, fill=ftlaranja!22] coordinates {\ftpulsos};
\nextgroupplot[title={(b) Núcleos de interpolação\\[1pt]{\tiny\ftstrut $x(nT_s)\,g(t-nT_s)/\tau$}}, clip=true]
  \ftnucleos
\nextgroupplot[title={(c) Soma dos núcleos\\[1pt]{\tiny\ftstrut $x_r'(t)/\tau$ e $x_r(t)$}}]
  \draw[black!45, densely dashed, thin] (axis cs:0,1) -- (axis cs:\fttmax,1);
  \addplot[gray!55, dashed, domain=0:\fttmax, samples=300] {cos(deg(2*pi*\ftfzero*x))};
  \addplot[ftroxo, thick, domain=0:\fttmax, samples=300, visible on=<2->]
    {\ftA*cos(deg(2*pi*\ftfzero*(x-\fttau/2)))};
  \addplot[ftverde, thick, densely dashed, domain=0:\fttmax, samples=300, visible on=<3->]
    {cos(deg(2*pi*\ftfzero*(x-\fttau/2)))};
\end{groupplot}
\end{tikzpicture}\par\vspace{2pt}
{\scriptsize
\tikz[baseline=-0.6ex]\draw[ftroxo, thick] (0,0) -- (0.4,0);\; após o passa-baixas
\qquad
\tikz[baseline=-0.6ex]\draw[ftverde, thick, densely dashed] (0,0) -- (0.4,0);\; após o equalizador
\qquad
\tikz[baseline=-0.6ex]\draw[gray!55, dashed] (0,0) -- (0.4,0);\; $x(t)$}
\end{center}
\pause
\end{frame}

% ============================================================================
%  Reconstrução no domínio da frequência (revelado em 3 etapas)
% ============================================================================
\begin{frame}{Método Topo-Plano: Reconstrução}{Representação e Demonstração no Domínio da Frequência}
\small
Na frequência, a reconstrução é o produto do espectro amostrado pelo passa-baixas ideal
e pelo equalizador, ambos definidos em $|f| < f_s/2$:
\begin{equation*}
  X_r(f) = X_p(f)\cdot \underbrace{T_s}_{H_{PB}(f)} \cdot\;
  \ftverde{\underbrace{\frac{1}{\tau\,\mathrm{sinc}(f\tau)}}_{H_{eq}(f)}}
  = e^{-j\pi f\tau}\,X(f)
  \quad\overset{\mathcal{F}^{-1}}{\longleftrightarrow}\quad
  \boxed{\,x_r(t) = x\!\left(t - \frac{\tau}{2}\right)}
\end{equation*}
O equalizador existe porque $\mathrm{sinc}(f\tau) \neq 0$ na banda: o primeiro nulo,
$1/\tau \geq f_s$, fica fora de $|f| < f_s/2$.
\begin{center}
\begin{tikzpicture}
\begin{groupplot}[group style={group name=g, group size=3 by 1, horizontal sep=1.3cm}, ftdemof, height=0.18\textheight]
\nextgroupplot[title={(a) Espectro amostrado\\[1pt]{\tiny\ftstrut $|X_p(f)|$ e $H_{PB}(f)$}}, yticklabels={$\frac{\tau}{2T_s}$}]
  \draw[ftroxo, densely dashed, fill=ftroxo!8] (axis cs:-\ftfc,0) rectangle (axis cs:\ftfc,1.15);
  \addplot[black!55, densely dashed, domain=-\ftfmax:\ftfmax, samples=400]
    {abs(sin(deg(pi*x*\fttau))/(pi*x*\fttau))};
  \addplot[ftimpulso, ftlaranja!90!black] coordinates {\ftrepsinc};
\nextgroupplot[title={(b) Equalizador\\[1pt]{\tiny\ftstrut $|H_{eq}(f)|$}}, yticklabels={$\frac{1}{\tau}$},
  ymax=1.8, visible on=<2->]
  \addplot[ftverde, thick, domain=-\ftfc:\ftfc, samples=100, fill=ftverde!8, visible on=<2->]
    {1/(sin(deg(pi*x*\fttau+1e-6))/(pi*x*\fttau+1e-6))} \closedcycle;
\nextgroupplot[title={(c) Sinal reconstruído\\[1pt]{\tiny\ftstrut $|X_r(f)|$}}, yticklabels={$\frac{1}{2}$}]
  \addplot[ftimpulso, ftverde, visible on=<3->] coordinates {(-\ftfzero,1) (\ftfzero,1)};
\end{groupplot}
\ftop{1}{2}{\times}
\ftop{2}{3}{=}
\end{tikzpicture}
\end{center}
\pause\pause
\end{frame}

% ============================================================================
%  Visão geral do processo completo (tempo e frequência).
%  Fica DEPOIS da reconstrução: os painéis "após o passa-baixas" e "após o
%  equalizador" só fazem sentido com as duas etapas já apresentadas.
% ============================================================================
\begin{frame}{Método Topo-Plano}{Visão Geral no Domínio do Tempo}
\centering
\begin{tikzpicture}
\begin{groupplot}[
  group style={group size=4 by 1, horizontal sep=0.85cm},
  ftaxis
]

% (a) Sinal original
\nextgroupplot[title={(a) Sinal original\\[1pt]{\tiny\ftstrut $x(t)=\cos(2\pi\cdot 2t)$}}]
  \addplot[ftazul, thick, domain=0:\fttmax, samples=400]
    {cos(deg(2*pi*\ftfzero*x))};

% (b) Sinal amostrado (topo plano)
\nextgroupplot[title={(b) Sinal amostrado\\[1pt]{\tiny\ftstrut $x_p(t)$}}]
  \addplot[gray!55, dashed, domain=0:\fttmax, samples=400]
    {cos(deg(2*pi*\ftfzero*x))};
  \addplot[ftlaranja!90!black, semithick, fill=ftlaranja!22, area legend] coordinates {\ftpulsos};
  \addplot[ycomb, ftazul, thin, mark=*, mark size=1.4pt,
    mark options={fill=ftazul, draw=white, line width=0.3pt}] coordinates {\ftamostras};

% (c) Após o passa-baixas: abertura + deslocamento
\nextgroupplot[title={(c) Após o passa-baixas\\[1pt]{\tiny\ftstrut $x_r'(t)$}},
  clip mode=individual]
  \draw[black!45, densely dashed, thin] (axis cs:0,1) -- (axis cs:\fttmax,1);
  \addplot[gray!55, dashed, domain=0:\fttmax, samples=400]
    {cos(deg(2*pi*\ftfzero*x))};
  \addplot[ftroxo, thick, domain=0:\fttmax, samples=400]
    {\ftA*cos(deg(2*pi*\ftfzero*(x-\fttau/2)))};
  \draw[gray, densely dotted] (axis cs:\ftpico,1) -- (axis cs:\ftpico,1.14);
  \draw[ftroxo, densely dotted] (axis cs:\ftpicoatr,\ftA) -- (axis cs:\ftpicoatr,1.14);
  \node[font=\tiny, anchor=south west, inner sep=1pt] at (axis cs:\ftpicoatr,1.12) {atraso $\tau/2$};

% (d) Após o equalizador: amplitude recuperada, atraso mantido
\nextgroupplot[title={(d) Após o equalizador\\[1pt]{\tiny\ftstrut $x_r(t)$}},
  clip mode=individual]
  \draw[black!45, densely dashed, thin] (axis cs:0,1) -- (axis cs:\fttmax,1);
  \addplot[gray!55, dashed, domain=0:\fttmax, samples=400]
    {cos(deg(2*pi*\ftfzero*x))};
  \addplot[ftverde, thick, domain=0:\fttmax, samples=400]
    {cos(deg(2*pi*\ftfzero*(x-\fttau/2)))};
  \draw[gray, densely dotted] (axis cs:\ftpico,1) -- (axis cs:\ftpico,1.14);
  \draw[ftverde, densely dotted] (axis cs:\ftpicoatr,1) -- (axis cs:\ftpicoatr,1.14);
  \node[font=\tiny, anchor=south west, inner sep=1pt] at (axis cs:\ftpicoatr,1.12) {atraso $\tau/2$};

\end{groupplot}
\end{tikzpicture}

\vspace{0.2em}
\mbox{\scriptsize
\tikz[baseline=-0.6ex]\draw[gray!55, dashed] (0,0) -- (0.45,0); $x(t)$\hspace{0.7em}
\tikz[baseline=-0.6ex]\draw[ftlaranja!90!black, semithick, fill=ftlaranja!22] (0,-0.08) rectangle (0.5,0.12);\hspace{3pt}pulsos de largura $\tau$\hspace{0.7em}
\tikz[baseline=-0.6ex]\fill[ftazul] (0,0) circle (1.4pt); \,amostras $x(nT_s)$\hspace{0.7em}
\tikz[baseline=-0.6ex]\draw[ftroxo, thick] (0,0) -- (0.45,0); após o passa-baixas\hspace{0.7em}
\tikz[baseline=-0.6ex]\draw[ftverde, thick] (0,0) -- (0.45,0); após o equalizador}
\end{frame}

\begin{frame}{Método Topo-Plano}{Visão Geral no Domínio da Frequência}
\centering
\begin{tikzpicture}
\begin{groupplot}[
  group style={group size=4 by 1, horizontal sep=0.85cm},
  ftfreqaxis
]

% (a) Espectro do sinal original
\nextgroupplot[title={(a) Espectro original\\[1pt]{\tiny\ftstrut $X(f)=\frac{1}{2}[\delta(f-2)+\delta(f+2)]$}}]
  \addplot[ftimpulso, ftazul] coordinates {(-\ftfzero,0.5) (\ftfzero,0.5)};

% (b) Espectro amostrado: réplicas em n f_s ponderadas pelo sinc
\nextgroupplot[title={(b) Sinal amostrado\\[1pt]{\tiny\ftstrut $X_p(f)$}}]
  \addplot[black!55, densely dashed, domain=-\ftfmax:\ftfmax, samples=500]
    {0.5*abs(sin(deg(pi*x*\fttau))/(pi*x*\fttau))};
  \addplot[ftimpulso, ftlaranja!90!black] coordinates {\ftreplicas};
  \addplot[ftimpulso, ftazul] coordinates {(-\ftfzero,\ftAmeio) (\ftfzero,\ftAmeio)};

% (c) Após o passa-baixas: só a banda base, atenuada pelo sinc
\nextgroupplot[title={(c) Após o passa-baixas\\[1pt]{\tiny\ftstrut $X_r'(f)$}}]
  \draw[black!45, densely dashed, thin] (axis cs:-\ftfmax,0.5) -- (axis cs:\ftfmax,0.5);
  \draw[ftroxo, densely dashed, fill=ftroxo!8] (axis cs:-\ftfc,0) rectangle (axis cs:\ftfc,0.6);
  \node[font=\tiny, anchor=south, inner sep=1pt, yshift=4pt, ftroxo] at (axis cs:0,0.6) {$H_{PB}(f)$};
  \addplot[ftimpulso, ftroxo] coordinates {(-\ftfzero,\ftAmeio) (\ftfzero,\ftAmeio)};

% (d) Após o equalizador: amplitude recuperada, fase linear mantida
\nextgroupplot[title={(d) Após o equalizador\\[1pt]{\tiny\ftstrut $X_r(f)$}}]
  \draw[black!45, densely dashed, thin] (axis cs:-\ftfmax,0.5) -- (axis cs:\ftfmax,0.5);
  \addplot[ftverde, densely dashed, domain=-\ftfc:\ftfc, samples=100, fill=ftverde!8]
    {0.5/(sin(deg(pi*x*\fttau+1e-6))/(pi*x*\fttau+1e-6))} \closedcycle;
  \node[font=\tiny, anchor=south, inner sep=1pt, yshift=4pt, ftverde!70!black] at (axis cs:0,0.6) {$H_{eq}(f)\propto\frac{1}{\mathrm{sinc}(f\tau)}$};
  \addplot[ftimpulso, ftverde] coordinates {(-\ftfzero,0.5) (\ftfzero,0.5)};

\end{groupplot}
\end{tikzpicture}

\vspace{0.2em}
\mbox{\scriptsize
\tikz[baseline=-0.6ex]\draw[black!55, densely dashed] (0,0) -- (0.45,0); envelope $\frac{1}{2}|\mathrm{sinc}(f\tau)|$ \hspace{0.7em}
\tikz[baseline=-0.6ex]\draw[ftazul, thick] (0,0) -- (0.45,0); banda base\hspace{0.7em}
\tikz[baseline=-0.6ex]\draw[ftlaranja!90!black, thick] (0,0) -- (0.45,0); réplicas em $nf_s \pm 2$}

\vspace{2pt}
\mbox{\scriptsize
\tikz[baseline=-0.6ex]\draw[ftroxo, densely dashed, fill=ftroxo!8] (0,-0.08) rectangle (0.5,0.12);\hspace{3pt}passa-baixas\hspace{0.7em}
\tikz[baseline=-0.6ex]\draw[ftverde, densely dashed, fill=ftverde!8] (0,-0.08) rectangle (0.5,0.12);\hspace{3pt}equalizador}
\end{frame}

%simulacao
% ============================================================================
%  Parâmetros do estudo — vêm ANTES da fundamentação teórica, porque todos os
%  métodos (ideal, natural e topo-plano) são demonstrados com estes valores.
% ============================================================================

\section{Parâmetros do Estudo}

\begin{frame}[allowframebreaks]{Parâmetros fixos}
    \small

    \textbf{Sinal original e espectro:}
    \begin{align*}
        x(t) &= \cos(2\pi \cdot 2t) + 0,5\cos(2\pi \cdot 18t) \\
        &\Updownarrow \\
        X(f) &= \frac{1}{2}[\delta(f-2) + \delta(f+2)]
        + \frac{0,5}{2}[\delta(f-18) + \delta(f+18)]
    \end{align*}

    \textbf{Filtro Anti-aliasing ($H_{AA}$):}
    \begin{align*}
        X_f(f) &= X(f) \cdot H_{AA}(f)
        \longrightarrow
        H_{AA}(f) =
        \begin{cases}
            1, & |f| \le 8\,\text{Hz}\\
            0, & |f| > 8\,\text{Hz}
        \end{cases}
        \quad \leftarrow f_c \\[2pt]
        X_f(f) &= \frac{1}{2}[\delta(f-2) + \delta(f+2)]
        \longleftrightarrow
        x_f(t) = \cos(2\pi \cdot 2t)
    \end{align*}

    \textbf{Parâmetros de Amostragem:}
    \[
        f_{max} = 2\,\text{Hz}
        \implies
        f_s > 2f_{max} = 4\,\text{Hz}
    \]

    \[
        f_s = 5\,\text{Hz}
        \implies
        T_s = \frac{1}{5}\,\text{s}
        \implies
        \text{banda de guarda} = 5 - 4 = 1\,\text{Hz}
    \]

    A partir dessa preparação, utiliza-se a equação filtrada em banda para realizar as amostragens: \begin{align*} x_f(t) = \cos(2\pi \cdot 2t) \end{align*}

\end{frame}

\begin{frame}{Parâmetros fixos}{Resumo}
    \small
    Estes são os valores usados em \textbf{todas} as demonstrações e figuras a seguir:
    \begin{center}
    \begin{tabular}{ll}
        \hline
        Sinal amostrado (pós anti-aliasing) & $x_f(t) = \cos(2\pi \cdot 2t)$ \\
        Frequência do sinal & $f_0 = f_{max} = 2$~Hz \\
        Corte do anti-aliasing & $f_c = 8$~Hz \\
        Frequência de amostragem & $f_s = 5$~Hz \\
        Período de amostragem & $T_s = 0{,}2$~s \\
        Banda de guarda & $f_s - 2f_{max} = 1$~Hz \\
        \hline
    \end{tabular}
    \end{center}

    \vspace{0.5em}
    Para cada método, o roteiro é o mesmo: descrição analítica geral, aplicação a
    $x_f(t)$ com estes parâmetros e, por fim, o filtro de reconstrução.
\end{frame}


\begin{frame}[allowframebreaks]{Referências}

	\bibliography{bibliography}

\end{frame}

\begin{frame}[plain]
	\vfill
	\centering
	{\color{themeblue2}\fontsize{40}{48}\selectfont\bfseries Obrigado!}
	\vfill
\end{frame}

\end{document}
